\documentclass[10pt,aspectratio=169,spanish]{beamer}

% --- PAQUETES Y CONFIGURACIÓN ---
\usepackage[utf8]{inputenc}
\usepackage[T1]{fontenc}
\usepackage[spanish,es-tabla]{babel}
\usepackage{amsmath,amssymb,amsfonts,amsthm}
\usepackage{mathtools}
\usepackage{tikz}
\usetikzlibrary{shapes.geometric}

% --- TEMA Y ESTILO VISUAL ---
\usetheme{Madrid}
\usecolortheme{whale}

% Colores Institucionales UFRO / Matemáticas
\definecolor{ufroblue}{RGB}{0, 51, 102}
\definecolor{ufrogold}{RGB}{212, 160, 23}
\definecolor{mathred}{RGB}{180, 40, 40}
\definecolor{mathgreen}{RGB}{30, 130, 60}

\setbeamercolor{palette primary}{bg=ufroblue,fg=white}
\setbeamercolor{palette secondary}{bg=ufroblue!85,fg=white}
\setbeamercolor{palette tertiary}{bg=ufroblue!70,fg=white}
\setbeamercolor{titlelike}{bg=ufroblue,fg=white}
\setbeamercolor{structure}{fg=ufroblue}
\setbeamercolor{block title}{bg=ufroblue,fg=white}
\setbeamercolor{block body}{bg=ufroblue!5,fg=black}
\setbeamercolor{block title example}{bg=mathgreen,fg=white}
\setbeamercolor{block body example}{bg=mathgreen!5,fg=black}
\setbeamercolor{block title alerted}{bg=mathred,fg=white}
\setbeamercolor{block body alerted}{bg=mathred!5,fg=black}

\setbeamertemplate{navigation symbols}{}
\setbeamertemplate{footline}[frame number]

% --- DATOS DEL CURSO Y DOCENTE ---
\title[Rectas y Circunferencias]{Geometría Analítica: La Recta y la Circunferencia}
\subtitle{Asignatura: Trigonometría y Geometría Analítica (IME078)}
\author[Prof. José Labrín Parra]{Prof. José  Labrín Parra}
\institute[UFRO]{
  Pedagogía en Matemática\\
  Departamento de Matemática y Estadística\\
  Facultad de Eduación y Humanidades\\
  Universidad de La Frontera
}
\date{Semestre II -- 2026}

\begin{document}

% -----------------------------------------------------------------------------
% PORTADA
% -----------------------------------------------------------------------------
\begin{frame}
  \titlepage
\end{frame}

% -----------------------------------------------------------------------------
% TEMARIO
% -----------------------------------------------------------------------------
\begin{frame}{Contenido Programático}
  \tableofcontents
\end{frame}

% -----------------------------------------------------------------------------
\section{Distancia entre Puntos y División de Segmentos}
% -----------------------------------------------------------------------------

\begin{frame}{1.1 Distancia entre dos puntos en el plano}
\begin{columns}
  \column{0.55\textwidth}
  \begin{block}{Teorema (Distancia Euclidiana)}
    Dados dos puntos $A(x_1, y_1)$ y $B(x_2, y_2)$ en el plano cartesiano $\mathbb{R}^2$, la distancia $d(A,B)$ está dada por:
    \[
      d(A,B) = \sqrt{(x_2 - x_1)^2 + (y_2 - y_1)^2}
    \]
  \end{block}
  \begin{proof}[Demostración corta]
    Construyendo el triángulo rectángulo $ABC$ con $C(x_2, y_1)$, se tiene $AC = |x_2 - x_1|$ y $BC = |y_2 - y_1|$. Por Teorema de Pitágoras:
    $d(A,B)^2 = AC^2 + BC^2$.
  \end{proof}

  \column{0.45\textwidth}
  \begin{center}
  \begin{tikzpicture}[scale=0.8]
    \draw[->,thick] (-0.5,0) -- (4.5,0) node[right] {$x$};
    \draw[->,thick] (0,-0.5) -- (0,4.5) node[above] {$y$};
    \coordinate (A) at (1,1);
    \coordinate (B) at (4,3.5);
    \coordinate (C) at (4,1);
    
    \draw[line width=1.2pt,ufroblue] (A) -- (B) node[midway,above left] {$d(A,B)$};
    \draw[dashed,gray] (A) -- (C) node[midway,below] {$|x_2-x_1|$};
    \draw[dashed,gray] (C) -- (B) node[midway,right] {$|y_2-y_1|$};
    
    \filldraw[mathred] (A) circle (2pt) node[below left] {$A(x_1,y_1)$};
    \filldraw[mathred] (B) circle (2pt) node[above right] {$B(x_2,y_2)$};
    \filldraw[gray] (C) circle (1.5pt) node[below right] {$C(x_2,y_1)$};
    \draw (3.7,1) -- (3.7,1.3) -- (4,1.3);
  \end{tikzpicture}
  \end{center}
\end{columns}
\end{frame}

\begin{frame}{1.2 División de un segmento en una razón dada}
\begin{block}{Punto de División $P(x,y)$}
Si $P_1(x_1,y_1)$ y $P_2(x_2,y_2)$ son extremos de un segmento, el punto $P(x,y)$ tal que $\frac{\vec{P_1P}}{\vec{PP_2}} = r$ ($r \neq -1$) posee coordenadas:
\[
  x = \dfrac{x_1 + r x_2}{1 + r}, \qquad y = \dfrac{y_1 + r y_2}{1 + r}
\]
\end{block}

\begin{columns}
\column{0.6\textwidth}
  \begin{alertblock}{Casos Especiales}
  \begin{itemize}
    \item \textbf{Punto Medio ($r=1$):} $M = \left(\frac{x_1+x_2}{2}, \frac{y_1+y_2}{2}\right)$.
    \item \textbf{Razón Negativa ($r < 0$):} El punto $P$ se ubica exterior al segmento $P_1P_2$.
  \end{itemize}
  \end{alertblock}
\column{0.4\textwidth}
  \begin{center}
  \begin{tikzpicture}[scale=0.7]
    \draw[->] (-0.5,0) -- (4.5,0);
    \draw[->] (0,-0.5) -- (0,3.5);
    \coordinate (P1) at (0.5,0.5);
    \coordinate (P) at (2.3,1.7);
    \coordinate (P2) at (3.5,2.5);
    
    \draw[line width=1pt,ufroblue] (P1) -- (P2);
    \filldraw (P1) circle (2pt) node[below] {$P_1$};
    \filldraw[mathred] (P) circle (2.5pt) node[above left] {$P(x,y)$};
    \filldraw (P2) circle (2pt) node[above] {$P_2$};
  \end{tikzpicture}
  \end{center}
\end{columns}
\end{frame}

\begin{frame}{1.3 Problemas Propuestos}
\begin{block}{Ejercicios Propuestos}
\begin{enumerate}
    \item \textbf{Distancia de Puntos:} Dados los puntos $A(-3, 4)$ y $B(5, -2)$, determine las coordenadas de un punto $C$ sobre el eje $Y$ de modo que el triángulo $ABC$ sea isósceles con $AC = BC$.
    \item \textbf{División de Segmento:} Sean $P_1(-2, 5)$ y $P_2(6, -3)$ los extremos del segmento $P_1P_2$. Halle las coordenadas de los puntos $A$ y $B$ que trisecan al segmento (es decir, lo dividen en tres partes de igual longitud).

      \item \textbf{División de Segmento:} Dados los puntos
$A(1,2)$ y $B(5,6)$
encuentra las coordenadas del punto \(P\) que divide externamente al segmento \(\overline{AB}\) en la razón
$
\frac{AP}{PB} = -\frac{3}{2}.
$



    \item \textbf{Distancia:} Determine el valor de $k$ para que la distancia entre $P(k, 2)$ y $Q(3, -1)$ sea exactamente $5$ unidades.
\end{enumerate}
\end{block}
\end{frame}


\begin{frame}{Soluciones}
\begin{enumerate}
    \item \textbf{Distancia de Puntos:} El punto $C$ sobre el eje $Y$ que cumple $AC = BC$ es
    \[
    \boxed{C\left(0, -\frac{1}{3}\right)}.
    \]

    \item \textbf{División de Segmento (trisección):} Los puntos que trisecan al segmento $\overline{P_1P_2}$ son
    \[
    \boxed{A\left(\frac{2}{3}, \frac{7}{3}\right)} \quad \text{y} \quad \boxed{B\left(\frac{10}{3}, -\frac{1}{3}\right)}.
    \]

    \item \textbf{División de Segmento (razón negativa):} El punto que divide externamente al segmento $\overline{AB}$ en la razón $-\frac{3}{2}$ es
    \[
    \boxed{P(13, 14)}.
    \]

    \item \textbf{Punto Medio / Distancia:} Los valores de $k$ para los cuales la distancia entre $P(k,2)$ y $Q(3,-1)$ es $5$ unidades son
    \[
    \boxed{k = -1 \quad \text{o} \quad k = 7}.
    \]
\end{enumerate}
\end{frame}

% -----------------------------------------------------------------------------
\section{Lugar Geométrico y la Ecuación de la Recta}
% -----------------------------------------------------------------------------

\begin{frame}{2.1 Concepto de Lugar Geométrico}
\begin{block}{Definición}
Un \textbf{Lugar Geométrico} es el conjunto de todos los puntos $P(x,y)$ en el plano que satisfacen una propiedad o condición geométrica común dada por una ecuación $F(x,y)=0$.
\end{block}

\begin{exampleblock}{Ejemplo Desarrollado}
Encuentre el lugar geométrico de los puntos $P(x,y)$ que equidistan de $A(1,-2)$ y $B(5,4)$.
\pause
\begin{align*}
  d(P,A) &= d(P,B) \implies (x-1)^2 + (y+2)^2 = (x-5)^2 + (y-4)^2 \\
  x^2 - 2x + 1 + y^2 + 4y + 4 &= x^2 - 10x + 25 + y^2 - 8y + 16 \\
  8x + 12y - 36 &= 0 \implies \mathbf{2x + 3y - 9 = 0} \quad \text{(Simetral/Mediatriz)}
\end{align*}
\end{exampleblock}
\end{frame}
\begin{frame}{2.1 La Recta como Lugar Geométrico}
\begin{columns}
\column{0.55\textwidth}
  \begin{block}{Definición Formal}
  Una \textbf{recta} $L$ es el lugar geométrico de todos los puntos $P(x, y)$ en el plano cartesiano tales que el cociente entre las diferencias de sus coordenadas con respecto a un punto fijo $P_0(x_0, y_0)$ es constante.
  \end{block}

  \begin{block}{Concepto de Pendiente ($m$)}
  Dicha constante se denomina \textbf{pendiente} ($m$) y representa la razón de cambio de $y$ respecto a $x$:
  \[
    m = \frac{\Delta y}{\Delta x} = \frac{y - y_0}{x - x_0} \quad \text{con } x \neq x_0
  \]
  Geométricamente, la pendiente equivale a la tangente del ángulo de inclinación $\alpha$: $m = \tan(\alpha)$.
  \end{block}

\column{0.45\textwidth}
  \begin{center}
  \begin{tikzpicture}[scale=0.75]
    % Ejes cartesianos
    \draw[->, thick] (-1,0) -- (4,0) node[right] {$X$};
    \draw[->, thick] (0,-1) -- (0,4) node[above] {$Y$};
    \node[below left] at (0,0) {$O$};

    % Recta L
    \draw[line width=1.2pt, ufroblue] (-0.5,0.5) -- (3.5,3.7) node[right] {$L$};

    % Puntos P0 y P
    \filldraw[mathred] (1,1.7) circle (2pt) node[above left] {$P_0(x_0,y_0)$};
    \filldraw[mathred] (2.8,3.14) circle (2pt) node[above left] {$P(x,y)$};

    % Triángulo incremental
    \draw[dashed, gray] (1,1.7) -- (2.8,1.7) node[midway, below] {\small $\Delta x$};
    \draw[dashed, gray] (2.8,1.7) -- (2.8,3.14) node[midway, right] {\small $\Delta y$};

    % Ángulo alfa
    \draw (1.5,1.7) arc (0:38:0.5);
    \node at (1.8,1.9) {\small $\alpha$};
  \end{tikzpicture}
  \end{center}
\end{columns}
\end{frame}


\begin{frame}{2.2 Formas de la Ecuación de la Recta}

  \begin{block}{1. Forma Punto-Pendiente y Principal}
  \begin{itemize}
    \item \textbf{Punto-Pendiente:} $y - y_0 = m(x - x_0)$
    \item \textbf{Principal (Pendiente-Intersección):} $y = mx + b$
  \end{itemize}
  \end{block}


  \begin{block}{3. Forma General}
  Ecuación lineal de primer grado:
  \[
    Ax + By + C = 0 \quad (m = -A/B)
  \]
  \end{block}

\end{frame}


\begin{frame}{2.2 Formas de la Ecuación de la Recta}
\begin{columns}
\column{0.58\textwidth}
  \begin{block}{2. Forma Simétrica (o Canónica)}
  Si la recta intersecta a los ejes en $(a, 0)$ y $(0, b)$ con $a, b \neq 0$:
  \[
    \frac{x}{a} + \frac{y}{b} = 1
  \]
  \textit{Donde $a$ es la abscisa en el origen y $b$ es la ordenada en el origen.}
  \end{block}


\column{0.42\textwidth}
  \begin{center}
  \begin{tikzpicture}[scale=0.75]
    % Ejes cartesianos
    \draw[->, thick] (-0.8,0) -- (3.5,0) node[right] {$X$};
    \draw[->, thick] (0,-0.8) -- (0,3.5) node[above] {$Y$};
    \node[below left] at (0,0) {$O$};

    % Recta que pasa por (2,0) y (0,2.5)
    \draw[line width=1.2pt, ufroblue] (-0.4,3) -- (2.8,-1) node[right] {$L$};

    % Puntos de corte
    \filldraw[mathred] (2,0) circle (2pt) node[above right] {$(a, 0)$};
    \filldraw[mathred] (0,2.5) circle (2pt) node[left] {$(0, b)$};

    % Etiquetas de las distancias a y b
    \draw[dashed, mathred] (0,0) -- (2,0) node[midway, above] { };
    \draw[dashed, mathred] (0,0) -- (0,2.5) node[midway, right] { };
  \end{tikzpicture}
  \end{center}
\end{columns}
\end{frame}



\begin{frame}{2.3 Paralelismo y Perpendicularidad}
\begin{columns}
\column{0.52\textwidth}
  \begin{block}{Rectas Paralelas ($L_1 \parallel L_2$)}
  Dos rectas no verticales son paralelas si y solo si sus pendientes son iguales (mismo ángulo de inclinación):
  \[
    L_1 \parallel L_2 \iff m_1 = m_2
  \]
  \end{block}

  \begin{block}{Rectas Perpendiculares ($L_1 \perp L_2$)}
  Dos rectas no verticales son perpendiculares si y solo si el producto de sus pendientes es $-1$:
  \[
    L_1 \perp L_2 \iff m_1 \cdot m_2 = -1 \iff m_2 = -\frac{1}{m_1}
  \]
  \end{block}

\column{0.48\textwidth}
  \begin{center}
  \begin{tikzpicture}[scale=0.7]
    \draw[->] (-1,0) -- (3.5,0) node[right] {$X$};
    \draw[->] (0,-1) -- (0,3.5) node[above] {$Y$};

    % Paralelas
    \draw[line width=1pt, ufroblue] (-0.5,0) -- (2.5,3) node[right] {$L_1$};
    \draw[line width=1pt, ufroblue] (0.5,0) -- (3.5,3) node[right] {$L_2$};

    % Perpendicular
    \draw[line width=1pt, mathred] (-0.5,2.5) -- (2.5,-0.5) node[right] {$L_3$};
    
    % Marca de ángulo recto
    \draw (0.9,1.6) -- (1.05,1.75) -- (0.9,1.9);
  \end{tikzpicture}
  \end{center}
\end{columns}
\end{frame}


  \begin{frame}{2.3 Paralelismo y Perpendicularidad}
\begin{block}{Ejercicios Propuestos}
\begin{enumerate}
    \item \textbf{Ecuación con Relación de Intersecciones:} Hallar la ecuación de la recta que pasa por el punto $(2,3)$ y cuya abscisa en el origen es el doble que la ordenada en el origen.

    \item \textbf{Punto Equidistante:} Hallar un punto de la recta $3x + y + 4 = 0$ que equidista de los puntos $(-5, 6)$ y $(3, 2)$.

    \item \textbf{Recta Perpendicular:} Hallar la ecuación de la recta de abscisa en el origen $-3/7$ y que es perpendicular a la recta $3x + 4y - 10 = 0$.

    \item \textbf{Perpendicular por Intersección:} Hallar la ecuación de la perpendicular a la recta $2x + 7y - 3 = 0$ en su punto de intersección con $3x - 2y + 8 = 0$.
\end{enumerate}
\end{block}
\end{frame}





\begin{frame}{Soluciones}
\begin{enumerate}
    \item \textbf{Ecuación con Relación de Intersecciones:} 
    La ecuación de la recta que pasa por $(2,3)$ con $a = 2b$ es
    \[
    \boxed{x + 2y - 8 = 0}.
    \]

    \item \textbf{Punto Equidistante:} 
    El punto sobre la recta $3x + y + 4 = 0$ equidistante a $(-5,6)$ y $(3,2)$ es
    \[
    \boxed{P(-2, 2)}.
    \]

    \item \textbf{Recta Perpendicular:} 
    La ecuación de la recta perpendicular a $3x + 4y - 10 = 0$ con $a = -3/7$ es
    \[
    \boxed{28x - 21y + 12 = 0}.
    \]

    \item \textbf{Perpendicular por Intersección:} 
    La recta perpendicular a $2x + 7y - 3 = 0$ en su punto de intersección con $3x - 2y + 8 = 0$ es
    \[
    \boxed{7x - 2y + 16 = 0}.
    \]
\end{enumerate}
\end{frame}

% ==========================================
% SECCIÓN 2.4: ÁNGULO ENTRE DOS RECTAS
% ==========================================

\begin{frame}{2.4 Ángulo entre dos Rectas}
\begin{columns}
\column{0.55\textwidth}
  \begin{block}{Teorema}
  El ángulo $\alpha$ medido en \textbf{sentido antihorario} desde la recta $L_1$ (con pendiente $m_1$) hacia la recta $L_2$ (con pendiente $m_2$) satisface:
  \[
    \tan(\alpha) = \frac{m_2 - m_1}{1 + m_2 m_1} \quad (m_1 m_2 \neq -1)
  \]
  \end{block}

  \begin{exampleblock}{Demostración}
  Por geometría exterior en el triángulo, $\theta_2 = \alpha + \theta_1 \implies \alpha = \theta_2 - \theta_1$.
  \[
    \begin{aligned}
      \tan(\alpha) &= \tan(\theta_2 - \theta_1) \\
      &= \frac{\tan(\theta_2) - \tan(\theta_1)}{1 + \tan(\theta_2)\tan(\theta_1)} = \frac{m_2 - m_1}{1 + m_2 m_1}
    \end{aligned}
  \]
  \end{exampleblock}

\column{0.45\textwidth}
  \begin{center}
  \begin{tikzpicture}[scale=0.75]
    % Ejes
    \draw[->, thick] (-1,0) -- (4,0) node[right] {$X$};
    \draw[->, thick] (0,-1.5) -- (0,4) node[above] {$Y$};
    \node[above left] at (0,0) {$O$};

    % Recta L1
    \draw[line width=1.2pt] (-1,-1) -- (3,3) node[below right] {$L_1$};
    % Recta L2
    \draw[line width=1.2pt] (0.5,3.5) -- (3.5,-1) node[right] {$L_2$};

    % Ángulo theta_1
    \draw[->, thick] (0.6,0) arc (0:45:0.6);
    \node at (0.9,0.3) {\small $\theta_1$};

    % Ángulo theta_2
    \draw[->, thick] (3.3,0) arc (0:118:0.5);
    \node at (3,0.7) {\small $\theta_2$};

    % Ángulo alpha
    \draw[->, thick] (2.1,2.1) arc (45:120:0.6);
    \node at (1.8,2.6) {$\alpha$};
  \end{tikzpicture}
  \end{center}
\end{columns}
\end{frame}


% ==========================================
% SECCIÓN 2.6: PROBLEMAS PROPUESTOS Y SOLUCIONES
% ==========================================

\begin{frame}{2.4 Ángulo entre dos Rectas}
\begin{block}{Ejercicios Propuestos}
\begin{enumerate}
    \item \textbf{Cálculo de Coordenadas:} El ángulo formado por la recta que pasa por los puntos $(-4, 5)$ y $(3, y)$ con la que pasa por $(-2, 4)$ y $(9, 1)$ es de $135^\circ$. Hallar el valor de $y$.

    \item \textbf{Pendiente entre Rectas:} La recta $L_2$ forma un ángulo de $60^\circ$ con la recta $L_1$. Si la pendiente de $L_1$ es $1$, hallar la pendiente de $L_2$.

    \item \textbf{Pendiente Desconocida:} Hallar la pendiente de una recta que forma un ángulo de $45^\circ$ con la recta que pasa por los puntos de coordenadas $(2, -1)$ y $(5, 3)$.

    \item \textbf{Ecuación de la Recta:} Hallar la ecuación de la recta que pasa por el punto $(2, 5)$ y forma un ángulo de $45^\circ$ con la recta $x - 3y + 6 = 0$.
\end{enumerate}
\end{block}
\end{frame}


\begin{frame}{ Soluciones}
\begin{enumerate}
    \item \textbf{Cálculo de Coordenadas:}
    El valor de $y$ para obtener un ángulo de $135^\circ$ es
    \[
    \boxed{y = 9}.
    \]

    \item \textbf{Pendiente entre Rectas:}
    La pendiente de la recta $L_2$ es
    \[
    \boxed{m_2 = -(2 + \sqrt{3})}.
    \]

    \item \textbf{Pendiente Desconocida:}
    La pendiente buscada es
    \[
    \boxed{m_2 = -7}.
    \]

    \item \textbf{Ecuación de la Recta:}
    La ecuación de la recta correspondiente es
    \[
    \boxed{2x - y + 1 = 0}.
    \]
\end{enumerate}
\end{frame}--

% ==========================================
% SECCIÓN 2.5: DISTANCIA DE UN PUNTO A UNA RECTA
% ==========================================

\begin{frame}{2.5 Distancia de un Punto a una Recta}
\begin{columns}
\column{0.55\textwidth}
  \begin{block}{Teorema de la Distancia Perpendicular}
  La distancia mínima $d$ desde un punto $P_0(x_0, y_0)$ a la recta de ecuación general $L: A x + B y + C = 0$ viene dada por:
  \[
    d(P_0, L) = \frac{|A x_0 + B y_0 + C|}{\sqrt{A^2 + B^2}}
  \]
  \end{block}

  \begin{exampleblock}{Ejemplo de Aplicación}
  Calcular la distancia de $P(0,1)$ a $L: -x + 2y - 7 = 0$:
  \[
    d = \frac{|(-1)(0) + 2(1) - 7|}{\sqrt{(-1)^2 + 2^2}} = \frac{|-5|}{\sqrt{5}} = \sqrt{5}
  \]
  \end{exampleblock}

\column{0.45\textwidth}
  \begin{center}
  \begin{tikzpicture}[scale=0.75]
    \draw[->] (-1,0) -- (4,0) node[right] {$X$};
    \draw[->] (0,-1) -- (0,4) node[above] {$Y$};
    \draw[line width=1.2pt, ufroblue] (-0.5,3.5) -- (3.5,0.5) node[right] {$L$};
    \filldraw[mathred] (2.5,3) circle (2pt) node[above right] {$P_0(x_0,y_0)$};
    \draw[dashed, mathred, thick] (2.5,3) -- (1.5,2) node[midway,above left] {$d$};
    \draw (1.35,2.15) -- (1.5,2.3) -- (1.65,2.15);
  \end{tikzpicture}
  \end{center}
\end{columns}
\end{frame}




\begin{frame}{Ejercicios Resueltos: Distancia Punto-Recta}
\begin{block}{Ejercicios Propuestos}
\begin{enumerate}
    \item \textbf{Distancia de un Punto a una Recta:} Hallar la distancia del punto $P(-2, 3)$ a la recta de ecuación general $L: 8x - 15y + 34 = 0$.

    \item \textbf{Familia de Paralelas a Distancia Dada:} Hallar las ecuaciones de las rectas paralelas a $8x - 15y + 34 = 0$ que distan $3$ unidades del punto $Q(-2, 3)$.
\end{enumerate}
\end{block}
\end{frame}


\begin{frame}{Soluciones: Distancia Punto-Recta}
\begin{enumerate}
    \item \textbf{Distancia de un Punto a una Recta:}
    \[
   \boxed{d = \frac{27}{17}\text{ unidades}}.
    \]

    \item \textbf{Familia de Paralelas a Distancia Dada:}

    \[
    \boxed{8x - 15y + 112 = 0} \quad \text{y} \quad \boxed{8x - 15y + 10 = 0}.
    \]
\end{enumerate}
\end{frame}

\section{La Circunferencia}
% -----------------------------------------------------------------------------

\begin{frame}{3.1 Ecuaciones de la Circunferencia}
\begin{block}{Definición}
Lugar geométrico de los puntos $P(x,y)$ cuya distancia a un punto fijo $C(h,k)$ (centro) es una constante $r > 0$ (radio).
\end{block}

\begin{columns}
\column{0.48\textwidth}
  \begin{block}{Forma Ordinaria / Canónica}
  \[ (x - h)^2 + (y - k)^2 = r^2 \]
  Si el centro es $C(0,0)$:
  \[ x^2 + y^2 = r^2 \]
  \end{block}

\column{0.48\textwidth}
  \begin{block}{Forma General}
  \[ x^2 + y^2 + Dx + Ey + F = 0 \]
  donde:
  $h = -\frac{D}{2}, \quad k = -\frac{E}{2}, \quad r = \frac{1}{2}\sqrt{D^2 + E^2 - 4F}$
  \end{block}
\end{columns}

\pause
\begin{alertblock}{Análisis del Discriminante de Radio ($\Delta = D^2 + E^2 - 4F$)}
\begin{itemize}
  \item Si $\Delta > 0$: Circunferencia real.
  \item Si $\Delta = 0$: Un único punto $(h,k)$ (circunferencia degenerada).
  \item Si $\Delta < 0$: Lugar geométrico vacío (conjunto vacío en $\mathbb{R}^2$).
\end{itemize}
\end{alertblock}
\end{frame}

\begin{frame}{3.2 Posiciones Relativas: Recta y Circunferencia}
Dada $C: (x-h)^2 + (y-k)^2 = r^2$ y la recta $L: y = mx + n$. Sustituyendo $y$ obtenemos $ax^2 + bx + c = 0$.

\begin{columns}
\column{0.3\textwidth}
  \begin{block}{Exterior ($\Delta < 0$)}
  \begin{center}
  \begin{tikzpicture}[scale=0.5]
    \draw (0,0) circle (1.2);
    \draw[ufroblue,thick] (-1.8,1.8) -- (1.8,1.8);
  \end{tikzpicture}
  \end{center}
  $d(C,L) > r$\\ Sin intersección.
  \end{block}

\column{0.3\textwidth}
  \begin{block}{Tangente ($\Delta = 0$)}
  \begin{center}
  \begin{tikzpicture}[scale=0.5]
    \draw (0,0) circle (1.2);
    \draw[mathgreen,thick] (-1.8,1.2) -- (1.8,1.2);
    \filldraw[mathred] (0,1.2) circle (2pt);
  \end{tikzpicture}
  \end{center}
  $d(C,L) = r$\\ 1 punto de contacto.
  \end{block}

\column{0.3\textwidth}
  \begin{block}{Secante ($\Delta > 0$)}
  \begin{center}
  \begin{tikzpicture}[scale=0.5]
    \draw (0,0) circle (1.2);
    \draw[mathred,thick] (-1.8,0.5) -- (1.8,0.5);
    \filldraw (-1.1,0.5) circle (2pt);
    \filldraw (1.1,0.5) circle (2pt);
  \end{tikzpicture}
  \end{center}
  $d(C,L) < r$\\ 2 puntos de corte.
  \end{block}
\end{columns}
\end{frame}

\begin{frame}{3.3 Tangente a una Circunferencia desde un punto exterior}
\begin{exampleblock}{Problema  Resuelto}
Encontrar las rectas tangentes a $(x-2)^2 + (y-1)^2 = 5$ que pasan por $P(-3,6)$.
\end{exampleblock}
\pause
\textbf{Solución:}
1. Familia de rectas por $P(-3,6)$: $y - 6 = m(x + 3) \implies mx - y + (3m + 6) = 0$.
2. Distancia del centro $C(2,1)$ a la recta debe ser igual al radio $r = \sqrt{5}$:
\[
  \frac{|m(2) - 1 + 3m + 6|}{\sqrt{m^2 + 1}} = \sqrt{5} \implies \frac{|5m + 5|}{\sqrt{m^2 + 1}} = \sqrt{5}
\]
3. Elevando al cuadrado:
\[
  (5m + 5)^2 = 5(m^2 + 1) \implies 25m^2 + 50m + 25 = 5m^2 + 5 \implies 20m^2 + 50m + 20 = 0
\]
Dividiendo por 10: $2m^2 + 5m + 2 = 0 \implies (2m+1)(m+2) = 0 \implies \mathbf{m_1 = -1/2}, \mathbf{m_2 = -2}$.
Rectas tangentes: $L_1: x + 2y - 9 = 0$ \quad y \quad $L_2: 2x + y = 0$.
\end{frame}


\begin{frame}{3.4 Problemas Propuestos: La Circunferencia}
\begin{block}{Ejercicios Propuestos}
\begin{enumerate}
    \item \textbf{Ecuación General:} Muestre que la ecuación $x^2 + y^2 - 6x + 4y + 13 = 0$ representa un punto (circunferencia degenerada) e identifique sus coordenadas.
    \item \textbf{Tangencia:} Encuentre la ecuación de la circunferencia centrada en $C(-1, 4)$ y que es tangente a la recta $L: 3x - 4y + 4 = 0$.
    \item \textbf{Intersección de Recta y Círculo:} Determine para qué valores del parámetro $k$ la recta $y = x + k$ es secante, tangente o exterior a la circunferencia $x^2 + y^2 = 8$.
\end{enumerate}
\end{block}
\end{frame}
% -----------------------------------------------------------------------------
\section{Ejercicios Avanzados: PAES y Olimpiadas}
% -----------------------------------------------------------------------------

\begin{frame}{4.1 Problema Tipo PAES (Competencia Matemática 2 - M2)}
\begin{alertblock}{Ejercicio Tipo PAES M2}
Un barco navega en el plano cartesiano siguiendo una trayectoria lineal dada por $L: 3x - 4y + 12 = 0$. Una estación de faro de transmisión delimitada por la circunferencia de cobertura $(x-2)^2 + (y+1)^2 = 16$ desea saber en qué punto la embarcación estará a la mínima distancia del faro $C(2,-1)$ y si logra ingresar a la zona de cobertura.
\end{alertblock}

\pause
\textbf{Resolución:}
\begin{enumerate}
  \item Distancia del centro del faro $C(2,-1)$ a la trayectoria del barco:
  \[
    d = \frac{|3(2) - 4(-1) + 12|}{\sqrt{3^2 + (-4)^2}} = \frac{|6 + 4 + 12|}{5} = \frac{22}{5} = 4.4
  \]
  \item Como $r = \sqrt{16} = 4$ y $d = 4.4 > 4$, la distancia mínima al centro es $4.4$.
  \item La distancia mínima al **borde de cobertura** del faro es $d - r = 4.4 - 4 = 0.4$ unidades.
  \item \textbf{Conclusión PAES:} El barco **no** ingresa a la zona de cobertura del faro.
\end{enumerate}
\end{frame}

\begin{frame}{4.2 Problema Tipo Olimpiada Matemática}
\begin{block}{Desafío de Olimpiada (Problema de Apolonio)}
Demuestre que el lugar geométrico de los puntos $P(x,y)$ del plano tales que el cociente de sus distancias a dos puntos fijos $A(-a,0)$ y $B(b,0)$ es una constante $k \neq 1$, corresponde a una circunferencia (Circunferencia de Apolonio).
\end{block}

\pause
\begin{proof}[Demostración]
Sea $\frac{d(P,A)}{d(P,B)} = k \implies d(P,A)^2 = k^2 d(P,B)^2$. En coordenadas:
$
  (x+a)^2 + y^2 = k^2 \left( (x-b)^2 + y^2 \right)
$
Expandiendo y reagrupando términos:
$
  (1 - k^2)x^2 + (1 - k^2)y^2 + 2(a + k^2 b)x + (a^2 - k^2 b^2) = 0
$
Como $k \neq 1$, podemos dividir todo por $1 - k^2$:
\[
  x^2 + y^2 + \left(\frac{2(a + k^2 b)}{1 - k^2}\right)x + \frac{a^2 - k^2 b^2}{1 - k^2} = 0
\]
Al completar cuadrados, esta expresión toma la forma ordinaria $(x-h)^2 + y^2 = R^2$, lo cual representa una circunferencia centrada en el eje $X$.
\end{proof}
\end{frame}
\end{document}
% -----------------------------------------------------------------------------
\section{Actividad Práctica en GeoGebra}
% -----------------------------------------------------------------------------

\begin{frame}{5. Laboratorio Práctico con GeoGebra}
\begin{exampleblock}{Instrucciones para el Taller en Laboratorio}
Abran GeoGebra Suite Calculadora / Geometría e introduzcan la siguiente secuencia de comandos para verificar conjeturas sobre tangencia:
\end{exampleblock}

\begin{enumerate}
  \item Crear punto $C = (2, 1)$ y deslizador de radio $r = \text{EjecutarSlider}(1, 5)$.
  \item Definir la circunferencia: \texttt{c = Círculo(C, r)}.
  \item Crear un punto exterior $P = (-3, 6)$.
  \item Usar la herramienta/comando: \texttt{Tangentes(P, c)}.
  \item Crear los puntos de tangencia $A$ y $B$ mediante \texttt{Intersección(L_1, c)}.
  \item Comprobar que los ángulos $\angle PAC$ y $\angle PBC$ son estrictamente de $90^\circ$ (\texttt{Ángulo(P, A, C)}).
\end{enumerate}
\end{frame}

% -----------------------------------------------------------------------------
% CIERRE
% -----------------------------------------------------------------------------
\begin{frame}
  \begin{center}
    \Huge\bfseries\color{ufroblue} ¡Muchas Gracias! \\[1cm]
    \Large Preguntas y Discusión
  \end{center}
\end{frame}

\end{document}